% R911: gerado; alterar dados e gerador, não esta prancha.

\refstepcounter{figure}\label{fig:AR06-mental}
\noindent{\large\sffamily\bfseries\color{navy}AR06 · Mapa mental}\par\smallskip
\noindent\textbf{Memória metacognitiva, observação e leitura visual}\par
\noindent\textit{Como recuperar experiências sem transformar memória em verdade ou instrução?}\par\medskip
\begin{center}

\tikzset{
 atlasnode/.style={draw=steel,fill=blueBg,rounded corners=2mm,align=center,
  text width=4cm,inner sep=5pt,font=\fontsize{9.5}{11.8}\selectfont,text=textmain},
 atlasroot/.style={atlasnode,fill=navy,text=white,text width=2.8cm,font=\fontsize{10}{12}\selectfont\bfseries},
 atlasdecision/.style={diamond,aspect=2.1,draw=amber,fill=amberBg,align=center,
  text width=3cm,inner sep=2pt,font=\fontsize{9.5}{11.5}\selectfont},
 atlascall/.style={-Stealth,draw=steel,thick},
 atlasevent/.style={-Stealth,draw=violet,dashed,thick},
 atlasdata/.style={-Stealth,draw=teal,densely dotted,thick},
 atlasstructure/.style={draw=slate,thick},
 atlasboundary/.style={-Stealth,draw=amber,dash dot,thick},
 atlasedgelabel/.style={fill=papel,inner sep=1.3pt,font=\fontsize{9}{10}\selectfont}
}

\begin{tikzpicture}
\node[atlasroot] (root) at (0,0) {AR06\\Memória metacognitiva, observação e leitura visual};

\node[atlasnode] (branch0) at (-4.5,2.5) {\textbf{Preparação}\\Consulta local\\Lições anteriores\\Limite de registros};
\draw[atlasstructure] (root) -- (branch0);

\node[atlasnode] (branch1) at (4.5,2.5) {\textbf{Registro}\\Contexto\\Resultado observado\\Próxima ação};
\draw[atlasstructure] (root) -- (branch1);

\node[atlasnode] (branch2) at (-4.5,-2.5) {\textbf{Avaliação}\\Heurística\\Sem calibração\\Sem promoção automática};
\draw[atlasstructure] (root) -- (branch2);

\node[atlasnode] (branch3) at (4.5,-2.5) {\textbf{Leitura visual}\\Snapshot\\Eventos\\Simulação identificada};
\draw[atlasstructure] (root) -- (branch3);
\end{tikzpicture}
\end{center}
{\small\textbf{Figura \thefigure.} Mapa mental --- Memória metacognitiva, observação e leitura visual}\par\smallskip
\noindent\textbf{Como funciona e por quê.} Parta do centro e leia os quatro ramos. As linhas associam conceitos; não representam chamadas de funções.\par
\noindent\textbf{Justificativa.} prepare reúne memórias e lições em briefing delimitado. complete registra o resultado informado, o erro observado e a próxima ação, preservando a confiança. Uma busca posterior pode recuperar esse episódio. O snapshot visual lê artefatos locais. A sequência é explicativa, condicionada à rota escolhida, e não transforma observação em comprovação.\par
\noindent\textbf{Limite.} O agrupamento é didático; cada responsabilidade será localizada no diagrama e conferida no fluxo.\par\smallskip
\fontecore{mci/metacognitive_loop.py}{MetacognitiveLoop.prepare}
\fontecore{mci/metacognitive_loop.py}{MetacognitiveLoop.complete}
